% main.tex -- 初步排版（骨架，可直接 pdflatex 编译；\todo 处待填）
% 编译前需把 figures/*.svg 转成 PDF：  for f in figures/*.svg; do rsvg-convert -f pdf -o ${f%.svg}.pdf $f; done
\documentclass[11pt]{article}

\usepackage[margin=1in]{geometry}
\usepackage{amsmath,amssymb,amsthm}
\usepackage{booktabs}
\usepackage{mathtools}
\usepackage{graphicx}
\usepackage{hyperref}
\usepackage{xcolor}

\newtheorem{theorem}{Theorem}
\newtheorem{proposition}{Proposition}
\newtheorem{lemma}{Lemma}
\newtheorem{corollary}{Corollary}
\theoremstyle{remark}
\newtheorem{remark}{Remark}

\newcommand{\todo}[1]{\textcolor{red}{[TODO: #1]}}
\newcommand{\norm}[1]{\lVert #1 \rVert}
\newcommand{\abs}[1]{\lvert #1 \rvert}
\DeclareMathOperator*{\argmin}{arg\,min}

\title{Running SPICE Sparse LU Without a Vendor Sparse Solver:\\
Mixed-Precision Refinement and a Pattern-Bound Roofline\\[2pt]
\large on a Domestic Accelerator}

\author{\todo{authors}}
\date{\today}

\begin{document}
\maketitle

\begin{abstract}
Vendor sparse direct solvers are frequently absent on emerging accelerators. On
our target the sparse interface of the vendor solver library returns
\emph{NOT\_SUPPORTED} for every circuit matrix we tried---two library versions,
four devices, and a disassembly showing the entry point compiled to return the
error code. A SPICE-grade sparse LU must therefore be assembled from dense
primitives. We show that this is possible with mixed-precision factorization and
iterative refinement, and we delimit when it pays.

Our claims are about feasibility and boundaries, not throughput. First, a
factorization computed at a low precision can be made to converge, and the
requirement is a pivot condition: rounding a high-precision factorization
preserves every pivot to relative accuracy, whereas recomputing at low precision
confines the smallest pivot to the noise level. The two constructions therefore
have the same backward error and opposite refinement behaviour. Second, the
attainable residual is governed by the componentwise growth factor
$\rho=\norm{\abs{L}\abs{U}}/\norm{A}$, which is computable from the factors
before any solve and predicts the residual to within two orders of magnitude
across eleven orders of magnitude in $\rho$; on circuit matrices it is $\rho$,
not $\kappa$, that separates the tractable matrices from the hard ones. Third, a
pattern-bound roofline shows that on this device the workload is compute-bound by
two orders of magnitude and that FP32 and FP64 execute at the same rate
($29.03$ versus $28.79$\,TFLOP/s), so lowering the precision of the arithmetic
buys no speed; the only lever is the efficiency of the kernel itself.

We do not claim throughput. Our prototype reaches $9.3\%$ of a vendor GEMM on the
same device, and we state the remaining gap explicitly: $1.5\times$ blocking,
$6.9\times$ kernel engineering, at most $5.9\times$ ordering. None of the
algorithmic conclusions depends on that gap.
\end{abstract}

\section{Introduction}

Write the same matrix twice at the same unit roundoff $u_q$: once by
\emph{compressing} a high-precision factorization, once by \emph{recomputing} it.
The backward errors agree to within $25\%$. The induced preconditioners do not:
$\norm{AM^{-1}-I}$ is $27$ in one case and $3.9\times10^{7}$ in the other
(Table~\ref{tab:dichotomy}). The reason is a \emph{pivot floor}---recomputation
confines the smallest pivot to the noise level, compression preserves it---and
the consequence is that the two constructions call for \emph{opposite} refinement
schemes.

That observation is the subject of this paper. It matters because the usual way
to lower the precision of a linear solve---``store the factors in half
precision''---mixes the two constructions without distinguishing them, and the
literature analyses the resulting algorithm while treating the factors as an
opaque object of some precision $u_f$. We show that the \emph{construction} is an
independent axis, and that the two axes interact: the accuracy attainable at a
given precision depends on how the factors were made.

\paragraph{The four slots.}
A solve has four places where precision can be lowered: the factorization
($O(n^3)$), the storage of $L$ and $U$ ($O(n^2)$), the triangular substitutions
($O(n^2)$ per solve), and the residual that drives refinement ($O(n^2)$ per
iteration). Only the first carries the arithmetic, so only the first buys
throughput; the fourth sets the attainable accuracy. Everything else is
bookkeeping. This asymmetry is what makes the design space tractable, and it is
what the cost model of Section~\ref{sec:structure} formalizes.

\paragraph{Three quantities.}
The analysis is organized by
\[
  \kappa(A)=\norm{A}\norm{A^{-1}},\qquad
  \rho=\frac{\norm{\abs{L}\abs{U}}}{\norm{A}},\qquad
  \omega=\norm{\abs{A^{-1}}\abs{L}\abs{U}}\le\kappa\rho .
\]
$\kappa$ amplifies the \emph{forward} error, $\rho$ the \emph{backward} error, and
$\omega$ is the componentwise condition number that appears when both are
combined. The practical point is that $\rho$ is computable in $O(\mathrm{nnz})$
from the factors alone, before any solve, and it predicts the residual to within
two orders of magnitude on the matrices we test.

\paragraph{Contributions.}
\begin{enumerate}
  \item A pivot-floor lemma (Lemma~\ref{lem:floor}) and a dichotomy theorem
        (Theorem~\ref{thm:dichotomy}) identifying the construction as an
        independent design axis, together with the assumptions (A1)--(A2) under
        which the theorem holds and a counterexample showing the assumptions are
        not automatic (Section~\ref{sec:dominance}).
  \item A two-sided error analysis (Propositions~\ref{prop:residual}
        and~\ref{prop:forward}) in which the residual refinement rate is the
        \emph{same} quantity that governs GMRES convergence, so that ``slow
        GMRES'' and ``slow refinement'' are one phenomenon rather than two.
  \item A hardware-parameterized cost model whose optimum is a corner solution
        (Proposition~\ref{prop:corner}), and its reformulation as an
        \emph{implementation threshold}: the low-precision tier pays only if the
        kernel exceeds $s_{11}^*(\kappa,n)$, a target the designer controls
        (Section~\ref{par:impl}).
  \item An empirical study on thirteen circuit matrices, where nine exhibit
        $\rho=1$ exactly and the other three are identified as hard by $\rho$
        alone, with no trial factorizations (Table~\ref{tab:res}).
\end{enumerate}

\paragraph{Roadmap.}
Section~\ref{sec:structure} develops the structural analysis;
Section~\ref{sec:bounds} the error bounds; Section~\ref{sec:exp} the experiments,
including two hardware platforms. Code and data are available in the repository
accompanying this submission.

\section{Setup}

Let $A\in\mathbb{R}^{n\times n}$ be nonsingular, $PAQ=LU$ with partial pivoting,
$M=L U$, and $u_q=2^{-p}$ the unit roundoff at $p$ significant bits. We measure
three quantities:
\begin{align}
  \kappa(A) &= \norm{A}_\infty \norm{A^{-1}}_\infty, \\
  \rho(A,L,U) &= \norm{\abs{L}\abs{U}}_\infty / \norm{A}_\infty, \\
  \omega(A,L,U) &= \norm{\abs{A^{-1}}\abs{L}\abs{U}}_\infty \;\le\; \kappa\,\rho .
\end{align}
Standard Gaussian elimination gives $\abs{\Delta A}\le\gamma_n\abs{L}\abs{U}$ with
$\gamma_n=nu_q/(1-nu_q)$, hence $\norm{\Delta A}_\infty\le \rho\,u_q\norm{A}_\infty$;
$\rho$ therefore controls the \emph{backward} error, while $\kappa$ controls the
\emph{forward} error.

% 03-section3.tex -- §3 完整版（直接 \input 进 main.tex 替换原 §3）
% 需要 main.tex 加一行:  \usepackage{mathtools}
%
% 3.1 两个构造   3.2 引理3.1 因子扰动   3.3 引理3.2 主元地板
% 3.4 (A1) 主元支配性（正例/反例）      3.5 命题3.3 压缩
% 3.6 (A2) 量化地板 + 命题3.4 重算      3.7 定理3.5 二分性
% 3.8 适用范围与失效模式

\section{Structure: pivot floors and the dichotomy}\label{sec:structure}

Throughout, $\norm{\cdot}$ denotes the spectral norm,
$\kappa=\kappa_2(A)=\sigma_{\max}(A)/\sigma_{\min}(A)$, and
$\rho=\norm{\abs{L}\abs{U}}_2/\norm{A}_2$.  The experiments of
Section~\ref{sec:exp} report $\infty$-norm analogues; the two differ by at most a
factor $\sqrt n$ and we do not distinguish them.

\subsection{Two constructions at the same precision}

Let $u\in(0,1)$ be the unit roundoff and $\gamma_k=ku/(1-ku)$.

\begin{definition}[Factor constructions]\label{def:constructions}
Let $A\in\mathbb{R}^{n\times n}$ be nonsingular and let $PAQ=LU$ be its
factorization in exact arithmetic, with $L$ unit lower triangular and
$\abs{L}\le1$ entrywise.
\begin{enumerate}
  \item \emph{Recompute}: $M_r=\tilde L\tilde U$, from Gaussian elimination with
        partial pivoting carried out in precision $u$.
  \item \emph{Compress}: $M_c=\hat L\hat U$ with
        $\hat L=\operatorname{fl}_u(L)$, $\hat U=\operatorname{fl}_u(U)$.
\end{enumerate}
\end{definition}

The constructions differ only in \emph{where} rounding is committed: compression
rounds the finished factors once; recomputation rounds every intermediate update.

\subsection{Perturbation of the factors}

\begin{lemma}\label{lem:perturb}
With $E_r=A-M_r$ and $E_c=A-M_c$,
\begin{equation}
  \abs{E_r}\le\gamma_{n}\abs{\tilde L}\abs{\tilde U},
  \qquad
  \abs{E_c}\le(2u+3u^2)\abs{L}\abs{U},
\end{equation}
hence
\begin{equation}\label{eq:same-backward}
  \norm{E_r}\le(1+O(u))\,\gamma_n\,\rho\norm{A},
  \qquad
  \norm{E_c}\le(2u+O(u^2))\,\rho\norm{A}.
\end{equation}
\end{lemma}

\begin{proof}
The first bound is the standard componentwise backward error of Gaussian
elimination \cite{higham2002,ch2018}. For the second, write
$L=\hat L+\delta L$, $U=\hat U+\delta U$ with $\abs{\delta L}\le u\abs{L}$ and
$\abs{\delta U}\le u\abs{U}$ (round to nearest). Then
\[
  A=LU=(\hat L-\delta L)(\hat U-\delta U)
   =\hat L\hat U-\delta L\,\hat U-\hat L\,\delta U+\delta L\,\delta U ,
\]
so $E_c=\delta L\,\hat U+\hat L\,\delta U-\delta L\,\delta U$ and
$\abs{E_c}\le(2u+3u^2)\abs{L}\abs{U}$ after absorbing $(1+u)\le 1+O(u)$.
\end{proof}

\begin{remark}\label{rem:same}
Lemma~\ref{lem:perturb} makes precise the sense in which the two constructions
have the \emph{same} backward error: both are $O(u\rho\norm{A})$, and neither
bound involves $\kappa$. Any difference in behaviour must therefore come from the
\emph{direction} of $E$, not its magnitude.
\end{remark}

\subsection{The pivot floor}

\begin{lemma}[Pivot floor]\label{lem:floor}
The computed pivots of $M_r$ satisfy
\begin{equation}\label{eq:floor}
  \abs{\tilde u_{kk}-u_{kk}}\;\le\;\gamma_{k-1}\norm{U},
  \qquad k=1,\dots,n .
\end{equation}
In particular, if $\abs{u_{kk}}\le\gamma_{k-1}\norm{U}$ the computed pivot is of
the order of the noise level $\gamma_{k-1}\norm{U}$ and carries no information
about $u_{kk}$.
\end{lemma}

\begin{proof}
The $(k,k)$ entry after $k-1$ updates is
$u_{kk}=a_{kk}-\sum_{j<k}l_{kj}u_{jk}$. In precision $u$ each of the $k-1$ terms
contributes a rounding error bounded by $u\abs{l_{kj}}\abs{u_{jk}}\le u\norm{U}$,
and the partial sums are themselves rounded, giving \eqref{eq:floor}.
\end{proof}

\begin{remark}[The floor is a many-update effect]\label{rem:fma}
A single fused multiply-add rounds its exact result once, so a $2\times2$
elimination commits \emph{relative} error and exhibits no floor. The floor
requires $k\gg1$ updates whose intermediate magnitudes are large while the final
value is small. We verified this numerically: for the family
$U=\bigl(\begin{smallmatrix}1&t\\0&\varepsilon\end{smallmatrix}\bigr)$ with $n=2$,
compression and recomputation agree in every printed digit for $t$ up to
$10^{6}$, whereas the $n=48$ experiments of Figure~\ref{fig:inv} show saturation
at $\kappa\approx u^{-1}$.
\end{remark}

\subsection{Pivot dominance}\label{sec:dominance}

Lemma~\ref{lem:floor} bounds how \emph{small} a computed pivot can be---one
direction only. Bounding $\norm{M^{-1}}$ requires relating pivots to singular
values.

\begin{definition}[Pivot dominance, (A1)]\label{def:A1}
A family of factorizations satisfies (A1) with constant $c>0$ if
\begin{equation}
  \sigma_{\min}(LU)\;\ge\;c\,\min_{k}\abs{u_{kk}}
  \qquad\text{for every }LU\text{ in the family.}
\end{equation}
\end{definition}

\begin{proposition}[Sufficient condition]\label{prop:A1-suff}
Let $U=V\Lambda V^{-1}$. Then
\begin{equation}
  \sigma_{\min}(LU)\;\ge\;\frac{\min_k\abs{u_{kk}}}{\kappa_2(V)\,\norm{L^{-1}}_2},
\end{equation}
so (A1) holds with $c^{-1}=\kappa_2(V)\norm{L^{-1}}_2$.
\end{proposition}

\begin{proof}
$\sigma_{\min}(U)=\sigma_{\min}(V\Lambda V^{-1})\ge
\sigma_{\min}(V)\min_k\abs{u_{kk}}\sigma_{\min}(V^{-1})
=\min_k\abs{u_{kk}}/\kappa_2(V)$, and
$\sigma_{\min}(LU)\ge\sigma_{\min}(L)\sigma_{\min}(U)$ with
$\sigma_{\min}(L)=1/\norm{L^{-1}}_2$.
\end{proof}

\paragraph{Where (A1) holds.}
Normal triangular factors ($\kappa_2(V)=1$); $M$-matrix factorizations with
nonnegative factors, where $\abs{L}\abs{U}=\abs{A}$ and $\norm{L^{-1}}$ is
controlled by diagonal dominance---this covers the circuit matrices of
Section~\ref{sec:exp}, nine of which give $\rho=1$ exactly.

\paragraph{Where (A1) fails.}
For $U=\bigl(\begin{smallmatrix}1&t\\0&\varepsilon\end{smallmatrix}\bigr)$ with
$t\gg1$ the eigenvalues are $1,\varepsilon$ but
$\kappa_2(V)\asymp t$ and
\[
  \sigma_{\min}(U)=\frac{\varepsilon}{\sqrt{1+t^2}}
  \;\approx\;\frac{\min_k\abs{u_{kk}}}{t},
\]
so (A1) fails with $c\approx t^{-1}$.  Numerically the ratio
$\sigma_{\min}/\min_k\abs{u_{kk}}$ falls from $7.1\times10^{-1}$ at $t=1$ to
$1.0\times10^{-6}$ at $t=10^{6}$, matching the prediction $1/\sqrt{1+t^2}$.

\begin{remark}[The two failure modes coincide]\label{rem:coincide}
The off-diagonal magnitude $t$ that destroys (A1) also raises the noise level
$\gamma_n\norm{U}\approx\gamma_nt$ that creates the pivot floor. Hence when (A1)
fails badly, the pivots are themselves sub-floor and \emph{no} low-precision
factorization---compressed or recomputed---is usable. The dichotomy below is then
vacuous rather than false; the matrix simply requires higher precision. This is
the analytical counterpart of the experimental ``red zone''.
\end{remark}


\begin{figure}[t]
  \centering
  \includegraphics[width=\linewidth]{figures/fig7_compress_vs_recompute.pdf}
  \caption{$\norm{M^{-1}}$ for compressed versus recomputed factors at the same
  $u_q$. Compression (solid) tracks $\kappa$; recomputation (dotted) saturates at
  $\approx1/(u\norm{U})$.}
  \label{fig:inv}
\end{figure}

\begin{figure}[t]
  \centering
  \includegraphics[width=\linewidth]{figures/fig8_dev_vs_kappa.pdf}
  \caption{$\norm{AM^{-1}-I}$: compression diverges with $\kappa$, recomputation
  is capped at $O(1/u_q)$. This is the quantity that Theorem~\ref{thm:dichotomy}
  separates and Proposition~\ref{prop:recurrence} turns into a contraction rate.}
  \label{fig:dev}
\end{figure}

\subsection{Compressed factors preserve the inverse norm}

\begin{proposition}\label{prop:compress}
Let $M_c$ be the compressed factorization and suppose $2u\rho\kappa<1$. Then
\begin{equation}
  \norm{M_c^{-1}}\;\le\;\frac{\kappa}{\norm{A}}\cdot\frac{1}{1-2u\rho\kappa},
  \qquad
  \norm{I-M_c^{-1}A}\;\le\;\frac{2u\rho\kappa}{1-2u\rho\kappa}.
\end{equation}
\end{proposition}

\begin{proof}
By Lemma~\ref{lem:perturb},
$\norm{E_c}\le2u\rho\norm{A}=2u\rho\kappa\,\sigma_{\min}(A)$.
Weyl's inequality gives
$\sigma_{\min}(M_c)\ge\sigma_{\min}(A)-\norm{E_c}\ge\sigma_{\min}(A)(1-2u\rho\kappa)$,
which yields the first bound since $\sigma_{\min}(A)=\norm{A}/\kappa$. For the
second, $I-M_c^{-1}A=M_c^{-1}E_c$, so
\[
  \norm{M_c^{-1}E_c}\le\frac{\kappa}{\norm{A}(1-2u\rho\kappa)}\cdot2u\rho\norm{A}
  =\frac{2u\rho\kappa}{1-2u\rho\kappa}. \qedhere
\]
\end{proof}

\begin{remark}
Proposition~\ref{prop:compress} is \emph{unconditional}: it uses only Weyl.
Compression therefore preserves the small singular values exactly when
$u\rho\kappa\ll1$---the regime in which stationary refinement is expected to work.
\end{remark}

\subsection{Recomputed factors bound the preconditioned deviation}

Weyl gives the same lower bound for $M_r$, since $\norm{E_r}$ and
$\norm{E_c}$ are the same size. The distinguishing feature of recomputation is
the \emph{upper} side: the floor of Lemma~\ref{lem:floor} prevents
$\sigma_{\min}(M_r)$ from collapsing below the noise level.

\begin{definition}[Quantization floor, (A2)]\label{def:A2}
The recomputed factorization satisfies (A2) with constant $c'>0$ if
$\sigma_{\min}(M_r)\ge c'\,u\norm{U}$.
\end{definition}

\begin{remark}
(A2) is the analytical form of the saturation in Figure~\ref{fig:inv}:
$\norm{M_r^{-1}}$ stops growing at $\approx1/(u\norm{U})$. It is a
\emph{generic} statement---rounding errors do not conspire to annihilate the
smallest singular direction---rather than a theorem; adversarial alignments
violate it, and then (A1) fails too by Remark~\ref{rem:coincide}.
\end{remark}

\begin{proposition}\label{prop:recompute}
Under (A2),
\begin{equation}
  \norm{AM_r^{-1}-I}=\norm{E_rM_r^{-1}}
  \;\le\;\frac{\gamma_n}{u}\cdot\frac{\rho\norm{A}}{c'\norm{U}}
  \;\approx\;\frac{n\,\rho\,\norm{L}}{c'},
\end{equation}
which is independent of $\kappa$ and of $u$.
\end{proposition}

\begin{proof}
$AM_r^{-1}-I=(A-M_r)M_r^{-1}=E_rM_r^{-1}$, so
$\norm{E_rM_r^{-1}}\le\norm{E_r}\norm{M_r^{-1}}
\le(\gamma_n\rho\norm{A})(c'u\norm{U})^{-1}
=\frac{\gamma_n}{u}\cdot\frac{\rho\norm{A}}{c'\norm{U}}$.
Since $\norm{A}=\norm{LU}\le\norm{L}\norm{U}$ and
$\gamma_n/u=n/(1-nu)$, the stated form follows.
\end{proof}

\subsection{The dichotomy}

\begin{theorem}[Dichotomy]\label{thm:dichotomy}
Let $M_c$ and $M_r$ be the two constructions at the same $u$.
\begin{enumerate}
  \item If $4u\rho\kappa<1$ then $\norm{I-M_c^{-1}A}<1$: compression is a
        contraction for stationary refinement, with rate $\Theta(u\rho\kappa)$.
  \item Under (A2), $\norm{AM_r^{-1}-I}\le n\rho\norm{L}/c'$, a constant
        independent of $\kappa$ and of $u$: recomputation yields a bounded
        preconditioner, and GMRES on $AM_r^{-1}=I+F$ converges in a number of
        iterations governed by $\norm{F}$ rather than by $\kappa$.
  \item Conversely, under (A1),
        $\norm{I-M_r^{-1}A}\gtrsim\norm{L}/c$ whenever the smallest pivot is
        sub-floor; in that regime recomputation is unusable for stationary
        refinement at \emph{every} $\kappa$, while compression still satisfies
        (1) as long as $\kappa\lesssim1/(u\rho)$.
\end{enumerate}
\end{theorem}

\begin{proof}
(1) follows from Proposition~\ref{prop:compress}: the bound is $<1$ iff
$4u\rho\kappa<1$, and the rate is the bound itself.
(2) is Proposition~\ref{prop:recompute}.
(3): $M_r^{-1}A=I-M_r^{-1}E_r$. Under (A1) and a sub-floor pivot,
$\sigma_{\min}(M_r)\approx c\,u\norm{U}$, so
$\norm{M_r^{-1}}\approx1/(cu\norm{U})$ while
$\norm{E_r}\approx\gamma_n\rho\norm{A}\approx\gamma_n\norm{L}\norm{U}$; the
product is $\approx(\gamma_n/u)\norm{L}/c\approx n\norm{L}/c\gg1$ when
$\norm{L}\gg1$. Compression is unaffected because its pivots are preserved
relatively (Lemma~\ref{lem:perturb} and Proposition~\ref{prop:compress}).
\end{proof}

\begin{corollary}[The two feasible sets]\label{cor:sets}
For a fixed residual target, the two refinement schemes require structurally
different factors:
\[
  \text{stationary refinement: } \norm{I-M^{-1}A}<1
  \ \Longrightarrow\ \text{compressed factors, } \kappa\lesssim\frac{1}{u\rho};
\]
\[
  \text{GMRES-based refinement: } \norm{AM^{-1}-I}=O(1)
  \ \Longrightarrow\ \text{recomputed factors, no }\kappa\text{ restriction
  from the factorization.}
\]
\end{corollary}

\subsection{Scope, constants, and failure modes}

\paragraph{The constants are loose, the scaling is not.}
Proposition~\ref{prop:compress} places the stationary-refinement cliff at
$\kappa\approx1/(4u\rho)$; the measurements of Section~\ref{sec:exp} put it at
$\kappa\approx6/u$ for $\rho\approx1$, a factor $24$ more generous. This is
expected: the bound takes the worst case over error directions, whereas the
actual iteration matrix has a spectrum clustered near zero except in a few
directions.

\paragraph{How to check (A1) and (A2) after the fact.}
Both are properties of the \emph{computed} factorization and can be measured
cheaply: (A1) by comparing $\sigma_{\min}(M)$ with $\min_k\abs{u_{kk}}$, and
(A2) by comparing $\sigma_{\min}(M_r)$ with $u\norm{U}$. We recommend reporting
both ratios next to $\rho$.

\paragraph{Relation to earlier analyses.}
Carson and Higham \cite{ch2018} analyse the \emph{iteration} and obtain
$\kappa u_f\ll1$ for stationary refinement; Higham and Mary \cite{hm2022} give
the five-precision framework and the condition
$\kappa^2u_f^2(u_g+\kappa u_p)\ll1$ for GMRES-based refinement. Our result is
orthogonal: we hold the algorithm fixed and vary the \emph{construction} of the
factors, showing that the two conditions above are attained by structurally
different objects and that neither can be substituted for the other.


% 04-section4.tex -- §4 完整版（\input 进 main.tex 替换原 §4）
% 与 §3 共用记号；核心是把 "残差的收缩因子" 等同于 §3 的 ‖AM^{-1}-I‖。

\section{Bounds: residual versus forward error}\label{sec:bounds}

Section~\ref{sec:structure} separated the two constructions by the size of
$\norm{AM^{-1}-I}$. That same quantity turns out to be the \emph{contraction
factor of the residual}, which is what makes the analysis actionable: a single
number predicts both the cost of GMRES and the attainable residual.

\subsection{The residual of a factorized solve}

Let $M$ be a factorization used to solve $Ax=b$, and let $\hat x$ be the computed
solution. Write $\hat x=M^{-1}b+\delta$, where $\delta$ collects the rounding of
the triangular substitutions. Then, exactly,
\begin{equation}\label{eq:residual-identity}
  r \;=\; b-A\hat x \;=\; (M-A)\,\hat x \;-\; A\,\delta \;=\; E\,\hat x - A\,\delta .
\end{equation}
The two terms are the \emph{factorization} error and the \emph{solve} error; their
relative sizes decide which one controls the residual.

\begin{proposition}[Residual bound]\label{prop:residual}
With $E$ as in Lemma~\ref{lem:perturb} and $\norm{\delta}\le c_tu_t\kappa(M)\norm{\hat x}$
for the substitution,
\begin{equation}
  \frac{\norm{r}}{\norm{A}\norm{\hat x}}
  \;\le\; \gamma_n\,\rho \;+\; c_t\,u_t\,\kappa(M).
\end{equation}
For a factorization at $u_f$ and substitutions at $u_t\le u_r\ll u_f\kappa$, the
first term dominates and
\begin{equation}
  \frac{\norm{r}}{\norm{A}\norm{\hat x}} \;\lesssim\; \gamma_n\,\rho .
\end{equation}
\end{proposition}

\begin{proof}
Immediate from \eqref{eq:residual-identity}, Lemma~\ref{lem:perturb} and the
triangle inequality:
$\norm{E\hat x}\le\gamma_n\norm{\abs{L}\abs{U}}\norm{\hat x}=\gamma_n\rho\norm{A}\norm{\hat x}$,
$\norm{A\delta}\le c_tu_t\kappa(M)\norm{A}\norm{\hat x}$.
\end{proof}

\begin{remark}[The $\gamma_n$ is worst case]\label{rem:gamma}
$\gamma_n=nu/(1-nu)$ counts the worst case in which every rounding in the
$n-1$ updates of an entry conspires. In practice the observed ratio
$\mathrm{res}/(\rho u)$ is $O(1)$--$O(10^2)$ and \emph{does not grow with $n$}:
for the $n=682{,}712$ matrix of Table~\ref{tab:res} it is $0.45$. The mechanism
is that the rounding errors of the individual updates are relative to their own
partial sums and add incoherently, whereas $\abs{L}\abs{U}$ already
over-estimates the entries that matter. We therefore use $\mathrm{res}\approx\rho u$
as an \emph{operational} estimate and $\gamma_n\rho u$ as the certified bound.
\end{remark}

\subsection{Refinement is a contraction driven by the same quantity}

\begin{proposition}[Residual recurrence]\label{prop:recurrence}
Let $x_{k+1}=x_k+M^{-1}r_k$ be one step of iterative refinement, with $r_k$
evaluated in precision $u_r$. Then
\begin{equation}\label{eq:recurrence}
  r_{k+1} \;=\; E M^{-1} r_k \;-\; A\,\delta_{k+1},
  \qquad\text{so}\qquad
  \norm{r_{k+1}} \;\le\; \norm{AM^{-1}-I}\,\norm{r_k} \;+\; \norm{A}\norm{\delta_{k+1}} .
\end{equation}
Consequently $\norm{r_k}\to$ a floor of order $\rho u_r\norm{A}\norm{x}$ at the
geometric rate $\norm{AM^{-1}-I}$.
\end{proposition}

\begin{proof}
$r_{k+1}=b-Ax_{k+1}=b-Ax_k-AM^{-1}r_k=r_k-AM^{-1}r_k-A\delta_{k+1}
=(I-AM^{-1})r_k-A\delta_{k+1}$, and $I-AM^{-1}=-EM^{-1}$.
The floor is Proposition~\ref{prop:residual} applied at precision $u_r$.
\end{proof}

\begin{remark}[Why \S\ref{sec:structure} and \S\ref{sec:bounds} are the same statement]
The contraction factor in \eqref{eq:recurrence} is $\norm{AM^{-1}-I}$: the very
quantity that Theorem~\ref{thm:dichotomy} bounds to $O(1/u_q)$ for recomputed
factors and to $\Theta(\kappa u_q)$ for compressed ones. A large
$\norm{AM^{-1}-I}$ therefore has two symptoms that are usually reported
separately---slow GMRES and slow residual refinement---and they are one
phenomenon.
\end{remark}

\subsection{Forward error from the residual}

\begin{proposition}[Forward bound]\label{prop:forward}
If $r=b-A\hat x$ then $x-\hat x=-A^{-1}r$ and
\begin{equation}
  \frac{\norm{x-\hat x}}{\norm{x}}
  \;\le\; \kappa\,\frac{\norm{r}}{\norm{A}\norm{x}}
  \;\lesssim\; \kappa\,\rho\,u_r .
\end{equation}
\end{proposition}

\begin{proof}
$x-\hat x=A^{-1}r$, so $\norm{x-\hat x}\le\norm{A^{-1}}\norm{r}
=\kappa\norm{r}/\norm{A}$; divide by $\norm{x}$ and insert
Proposition~\ref{prop:residual} at precision $u_r$.
\end{proof}

\begin{corollary}[Two independent targets]\label{cor:targets}
A residual target $\tau$ (what a SPICE-like convergence check enforces) needs
$\rho u_f\lesssim\tau$ and is \emph{independent of $\kappa$};
a forward-error target $\varepsilon$ needs $\kappa\rho u_r\lesssim\varepsilon$
and is \emph{linear in $\kappa$}. The two requirements are decoupled by the
choice of $u_r$: one may stop at $u_f$ (cheap) when only the residual is
checked, and must refine to $u_r\approx\varepsilon/(\kappa\rho)$ otherwise.
\end{corollary}

\subsection{Empirical validation}

\begin{table}[t]
  \centering
  \small
  \caption{$\mathrm{res}/(\rho u_{64})$ on circuit matrices (BTF, AMD, KLU).
  Ratios stay within two orders of magnitude across eleven orders of magnitude
  in $\rho$; full table in the repository.}
  \label{tab:res}
  \begin{tabular}{lrrrrr}
    \toprule
    matrix & $n$ & $\rho$ & $\kappa_\infty$ & residual & $\mathrm{res}/(\rho u_{64})$ \\
    \midrule
    ASIC\_100k  & 99{,}340  & $1.000$ & $1.5\times10^{11}$ & $4.3\times10^{-17}$ & $0.39$ \\
    ASIC\_680ks & 682{,}712 & $1.000$ & $9.5\times10^{19}$ & $4.9\times10^{-17}$ & $0.45$ \\
    bcircuit    & 68{,}902  & $1.000$ & $2.2\times10^{9}$  & $7.3\times10^{-18}$ & $0.07$ \\
    G2\_circuit & 150{,}102 & $1.004$ & $2.0\times10^{7}$  & $7.7\times10^{-15}$ & $69$ \\
    rajat25     & 87{,}190  & $6.9\times10^{7}$ & $2.0\times10^{15}$ & $9.6\times10^{-9}$ & $1.25$ \\
    rajat27     & 20{,}640  & $2.1\times10^{4}$ & $2.0\times10^{15}$ & $1.8\times10^{-12}$ & $0.78$ \\
    twotone     & 120{,}750 & $1.6\times10^{11}$ & $4.5\times10^{9}$ & $1.4\times10^{-6}$ & $0.08$ \\
    \bottomrule
  \end{tabular}
\end{table}

Proposition~\ref{prop:residual} predicts $\mathrm{res}\lesssim\gamma_n\rho u$;
Table~\ref{tab:res} shows $\mathrm{res}\approx\rho u$ with a constant between
$0.07$ and $74$, and no systematic growth with $n$. Two consequences are worth
recording:
\begin{enumerate}
  \item The bound is usable as a \emph{planning} tool: $\rho$ is computable in
        $O(\mathrm{nnz})$ from the factors alone, before any solve is performed.
  \item The residual target of $10^{-6}$ used by circuit simulators is met by
        FP32 on the ten matrices with $\rho\approx1$ ($\rho u_{32}\approx6\times10^{-8}$)
        and missed by the three with $\rho\ge2\times10^{4}$---a split that
        $\kappa$ alone does not predict (ASIC\_680ks has $\kappa=9.5\times10^{19}$
        yet $\rho=1$).
\end{enumerate}

\subsection{Remarks}

\paragraph{Relation to componentwise backward error.}
Proposition~\ref{prop:residual} is the normwise consequence of the componentwise
statement $\abs{r}\le\gamma_n\abs{L}\abs{U}\abs{\hat x}$, which is the LU
analogue of the Oettli--Prager criterion. Reporting $\rho$ rather than
$\kappa$ is therefore not a heuristic: it is the natural scale of the backward
error of a factorized solve.

\paragraph{What the analysis does not cover.}
The bound assumes the substitutions are accurate for $M$; if the factors
themselves are so ill-conditioned that the solve loses digits ($\kappa(M)u_t\gtrsim1$,
the regime of Theorem~\ref{thm:dichotomy}(3)), the second term of
Proposition~\ref{prop:residual} dominates and the residual is limited by the
solve rather than by the factorization. This is exactly the red zone of
Remark~\ref{rem:coincide}.


\section{Hardware: the $(\kappa,s)$ phase diagram}

Let $s_u$ be the throughput of precision $u$ relative to FP64 ($s_{64}=1$, $s_{32}=s$,
$s_{11}=4s$), and $k(u,\kappa)$ the measured number of refinement steps. In
FP64-equivalent units
\begin{equation}
  T(u) = \frac{2}{3}\frac{n^3}{s_u} + 4k(u,\kappa)\,n^2 ,
  \qquad
  S(u) = \frac{\frac{2}{3}n^3}{T(u)} .
\end{equation}

\begin{proposition}[Corner solution and crossover]\label{prop:corner}
$T$ is monotone in every precision slot and the feasible set is monotone in each, so
the optimum lies on the boundary; for large $n$ it is a corner solution. Between two
precision levels the winner flips at
\begin{equation}
  n^*(\kappa,s) \;=\; \frac{6\,k(\kappa)\,s}{s-1}.
\end{equation}
\end{proposition}

\paragraph{The implementation axis.}\label{par:impl}
Proposition~\ref{prop:corner} treats $s_{32}$ and $s_{11}$ as hardware
constants. They are not: $s_{32}$ is fixed by the device, but $s_{11}$ is a
property of the \emph{kernel}. Writing $\Delta k=k_{11}-k_{32}$ for the extra
refinement steps the low-precision tier needs, that tier beats FP32 if and only
if
\begin{equation}
  s_{11} \;>\; s_{11}^*(\kappa,n) \;=\; \frac{1}{1/s_{32}-6\Delta k/n},
  \qquad\text{solvable only when } n>6\,s_{32}\,\Delta k .
\end{equation}
The threshold decreases monotonically in $n$ and approaches $s_{32}$ from above,
so the low-precision tier is never \emph{cheap} to enter---it is entered from
above.

Measured on the Mars X201, $s_{32}=2.288$ and a half-storage /
FP32-accumulate kernel reaches only $s_{16}=2.339$. Table~\ref{tab:s11} and
Figure~\ref{fig:s11} show that this is $2\%$--$17\%$ short of the threshold at
the sizes we can afford, so the tier does not pay in that implementation---but a
tensor-core path, whose $s_{11}$ is plausibly $8$--$16$, clears the threshold by
a factor of three or more. The decision variable is therefore the
\emph{implementation}, not the device: the same accelerator supports two
different optima depending on how the low-precision kernel is written.

\begin{table}[t]
  \centering
  \small
  \caption{Required $s_{11}$ to make the 11-bit tier pay ($s_{32}=2.288$).}
  \label{tab:s11}
  \begin{tabular}{lrrrr}
    \toprule
    $\kappa$ & $\Delta k$ & $n>6s_{32}\Delta k$ & $s_{11}^*(n=4096)$ & $s_{11}^*(n=16384)$ \\
    \midrule
    $10^{6}$  & 79  & 477 & 1.14 & 1.04 \\
    $10^{8}$  & 110 & 665 & 1.20 & 1.05 \\
    $10^{10}$ & 78  & 471 & 1.14 & 1.04 \\
    $10^{12}$ & 48  & 290 & 1.08 & 1.03 \\
    \bottomrule
  \end{tabular}
\end{table}

\begin{figure}[t]
  \centering
  \includegraphics[width=\linewidth]{figures/fig10_s11_threshold.pdf}
  \caption{The implementation threshold $s_{11}^*(\kappa,n)$. Curves: required
  low-precision throughput; horizontal lines: the measured non-MMA implementation
  ($2.34$) and a plausible tensor-core range ($\ge8$).}
  \label{fig:s11}
\end{figure}

\begin{figure}[t]
  \centering
  \includegraphics[width=\linewidth]{figures/fig9_phase_n1024.pdf}
  \caption{$(\kappa,s)$ phase diagram at $n=1024$. Colour marks the optimal precision;
  each hardware platform is a horizontal slice.}
  \label{fig:phase}
\end{figure}

% 05-section6.tex -- §6 实验（可投稿结构）
% 依赖：03/04/05 节定义的标签；图 fig1/fig2 在本节首次出现。

\section{Experiments}\label{sec:exp}

We ask three questions in order: (i) where are the phase boundaries of
Theorem~\ref{thm:dichotomy} located, (ii) does the pivot-floor mechanism explain
them, and (iii) does the resulting rule transfer to circuit-scale matrices and to
two different hardware stacks.

\subsection{Setup}

\paragraph{Precision emulator.}
All dense experiments use a pure-Python emulator that rounds \emph{every}
multiply-add to $p$ significant bits. Since $a\cdot b$ is exact in double for
$p\le27$, the operation $\mathrm{round}_p(a\cdot b+c)$ models a single-rounded
FMA, with no hidden extended precision. This gives direct control over the three
slots $(u_f,u_t,u_r)$ that a BLAS call cannot separate.

\paragraph{Matrix families.}
Condition-number-controlled matrices $A=U\Sigma V^{\top}$ with log-spaced
singular values (orthonormal factors from two-pass modified Gram--Schmidt), plus
1-D and 2-D Laplacians, diagonally dominant matrices, MNA-like sparse matrices,
and Wilkinson's growth example.

\paragraph{Protocol.}
$n=64$ unless stated; eleven variants $\times$ ten condition numbers
$\kappa\in\{1,10^{2},\dots,10^{14}\}$ $\times$ five seeds; spot checks at
$n=32,128$ and a banded $n=576$. The eleven variants differ only in
$(u_f,u_t,u_r)$ and in whether the factors are compressed or recomputed.

\paragraph{Circuit matrices.}
Thirteen matrices from the SuiteSparse circuit collection, factorized with KLU
under BTF with AMD within blocks. Their order ranges from $4{,}960$ to
$682{,}712$ and their sparsity from $0.08\%$ to $0.4\%$.

\paragraph{Metrics.}
\emph{Forward:} $\mathrm{FE}/\mathrm{FE}_{64}$, the ratio of the computed forward
error to that of a pure FP64 solve, so that a value near $1$ means ``as accurate
as FP64''. \emph{Backward:} the standard relative residual
$\norm{r}/(\norm{A}\norm{\hat x}+\norm{b})$. \emph{Cost:} FP64-equivalent flops
with $s_{32}=2$ and $s_{11}=8$ unless a measured value is quoted.

\subsection{Where the phase boundaries are}

\begin{table}[t]
  \centering
  \small
  \caption{Forward-error ratio $\mathrm{FE}/\mathrm{FE}_{64}$ (median of five
  seeds, $n=64$). Values near $1$ mean FP64 accuracy is recovered; values above
  $10^{5}$ mean the refinement diverged.}
  \label{tab:variants}
  \begin{tabular}{lrrrrrrr}
    \toprule
    variant & $10^{2}$ & $10^{4}$ & $10^{6}$ & $10^{8}$ & $10^{9}$ & $10^{10}$ & $10^{12}$ \\
    \midrule
    FP32, no refinement      & $2.3{\cdot}10^{8}$ & $2.5{\cdot}10^{8}$ & $1.4{\cdot}10^{8}$ & $1.7{\cdot}10^{8}$ & $4.9{\cdot}10^{8}$ & $5.4{\cdot}10^{7}$ & $2.0{\cdot}10^{6}$ \\
    FP32 + stationary        & $0.84$ & $1.0$ & $1.6$ & $2.0$ & $7.5{\cdot}10^{9}$ & $9.0{\cdot}10^{7}$ & $6.3{\cdot}10^{6}$ \\
    FP32 + GMRES             & $0.70$ & $1.1$ & $1.2$ & $0.69$ & $1.8$ & $1.2$ & $1.7$ \\
    11-bit recompute + IR    & $0.66$ & $1.5{\cdot}10^{11}$ & $1.7{\cdot}10^{12}$ & $3.1{\cdot}10^{10}$ & $2.1{\cdot}10^{11}$ & $2.3{\cdot}10^{8}$ & $2.7{\cdot}10^{6}$ \\
    11-bit recompute + GMRES & $0.86$ & $1.3$ & $1.6$ & $1.4$ & $1.4$ & $1.4$ & $1.7$ \\
    11-bit compress + IR     & $0.87$ & $0.92$ & $1.0$ & $7.3{\cdot}10^{9}$ & $3.7{\cdot}10^{12}$ & $9.6{\cdot}10^{13}$ & $9.1{\cdot}10^{18}$ \\
    11-bit compress + GMRES  & $0.72$ & $1.2$ & $1.2$ & $1.3$ & $1.2$ & $1.2$ & $1.3{\cdot}10^{7}$ \\
    \bottomrule
  \end{tabular}
\end{table}

Table~\ref{tab:variants} and Figure~\ref{fig:envelope} locate the boundaries.
Stationary refinement with FP32 factors collapses between $\kappa=10^{8}$ and
$10^{9}$, whereas the guarantee tabulated by Carson and Higham
\cite{ch2018,hm2022} is $\kappa\le10^{8}$: our experiment is tighter than
the theory, by construction, because the bound is worst case. GMRES-based
refinement with the same factors survives to $\kappa=10^{14}$, at which point the
FP64 forward error is itself $10^{-2}$ and nothing better is attainable.

\begin{figure}[t]
  \centering
  \includegraphics[width=\linewidth]{figures/fig1_envelope.pdf}
  \caption{The highest condition number at which each configuration still
  recovers FP64 accuracy (median over five seeds). Three configurations reach
  $10^{14}$: that is the FP64 floor, not the algorithm's limit.}
  \label{fig:envelope}
\end{figure}

\begin{figure}[t]
  \centering
  \includegraphics[width=\linewidth]{figures/fig2_cliff.pdf}
  \caption{Forward error versus $\kappa$. All configurations track the FP64
  floor $\kappa u_{64}$ (dashed) and then leave it vertically: the transition is
  a cliff, not a degradation.}
  \label{fig:cliff}
\end{figure}

The iteration counts grow approximately linearly in $\log\kappa$: for FP32
factors with GMRES refinement the median total is $7,13,37,53,83,107$ steps at
$\kappa=10^{6},10^{8},10^{9},10^{10},10^{12},10^{14}$. This growth is the price
of the extended range and it is what Proposition~\ref{prop:corner} converts into
a crossover in $n$.

\subsection{Does the pivot floor explain it}

\begin{table}[t]
  \centering
  \small
  \caption{Compressed versus recomputed factors at the same $u_q=2^{-11}$,
  $\kappa=10^{12}$, $n=48$. The backward errors agree; the preconditioner
  deviations differ by six orders of magnitude.}
  \label{tab:dichotomy}
  \begin{tabular}{lrrrr}
    \toprule
    construction & backward error & $\abs{\hat u_{nn}}$ & $\norm{M^{-1}}$ & $\norm{AM^{-1}-I}$ \\
    \midrule
    FP64 factors        & $1.9{\cdot}10^{-16}$ & $5.6{\cdot}10^{-11}$ & $1.9{\cdot}10^{11}$ & $4.5{\cdot}10^{-5}$ \\
    recompute at 11 bits & $3.0{\cdot}10^{-4}$  & $3.0{\cdot}10^{-5}$  & $2.2{\cdot}10^{5}$  & $27$ \\
    compress to 11 bits  & $3.8{\cdot}10^{-4}$  & $5.6{\cdot}10^{-11}$ & $1.9{\cdot}10^{11}$ & $3.9{\cdot}10^{7}$ \\
    \bottomrule
  \end{tabular}
\end{table}

Table~\ref{tab:dichotomy} is the mechanism check. The two low-precision
constructions have the same backward error to within $25\%$, yet the last pivot
is preserved by compression ($5.6\cdot10^{-11}$, matching FP64) and floored by
recomputation ($3.0\cdot10^{-5}\approx u_q\norm{U}$). Consequently
$\norm{AM^{-1}-I}$ is bounded at $27$ for recomputation and grows to
$3.9\cdot10^{7}$ for compression, exactly as Proposition~\ref{lem:floor} and
Theorem~\ref{thm:dichotomy} predict. This table is also the empirical content of
the ``same backward error, opposite fate'' claim in the abstract.

\subsection{Circuit-scale matrices}

Table~\ref{tab:res} reports $\rho$ and the residual on thirteen circuit
matrices. Two observations go beyond the dense sweep.

First, nine matrices have $\rho=1.000$ to three digits. This is not a
coincidence: for an $M$-matrix with nonnegative factors, $\abs{L}\abs{U}=\abs{A}$
exactly, so the componentwise growth is one. On such matrices a low-precision
factorization is componentwise backward stable regardless of $\kappa$; the
residual target of $10^{-6}$ is met by FP32 with two orders of margin
($\rho u_{32}\approx6\cdot10^{-8}$).

Second, the growth factor identifies the hard matrices without trial. The three
matrices with $\rho\ge2\cdot10^{4}$ are precisely those flagged as
near-singular by the solver and the slowest in our runs. This makes $\rho$ a
cheap a-priori tag: it is computable in $O(\mathrm{nnz})$ from the factors,
before any solve.

\paragraph{Ordering interacts with $\rho$.}
For the matrix \texttt{rajat25}, blocking the triangular form raises the fill-in
from $7.2$M to $42.5$M nonzeros and $\rho$ from $37.7$ to $6.9\cdot10^{7}$, a
change of six orders of magnitude on the same matrix. The heuristic that
switches blocking off when the predicted fill exceeds $20\times\mathrm{nnz}$ is
therefore not merely a performance rule: it also protects the numerical quality
of the factors.

\subsection{Two platforms, and the implementation axis}

\begin{table}[t]
  \centering
  \small
  \caption{Measured throughput ratio $s=\mathrm{TP}/\mathrm{TP}_{64}$ for a
  shared-memory GEMM microbenchmark (checksum-verified).}
  \label{tab:platforms}
  \begin{tabular}{llrr}
    \toprule
    platform & kernel & $s$ & note \\
    \midrule
    aarch64 CPU, 128 cores & blocked LU, FP32 & $1.88$--$2.46$ & median $2.10$ \\
    MetaX Mars X201 & GEMM, FP32 & $2.288$ & $n=8192$ \\
    MetaX Mars X201 & GEMM, half storage + FP32 accumulate & $2.339$ & no MMA \\
    \bottomrule
  \end{tabular}
\end{table}

Table~\ref{tab:platforms} shows that the two platforms agree to within $9\%$ on
$s_{32}$: for the FP32/FP64 pair they sit in the same regime, and the precision
decision is governed by $\kappa$ and $n$, not by the device.

Lowering the \emph{storage} precision does not change this: a half-storage kernel
that still accumulates in FP32 reaches $s_{16}=2.339$, only $2\%$ above
$s_{32}$, because the multiply-adds remain on the FP32 units. Without a
tensor-core path the low-precision regime of Figure~\ref{fig:phase} is therefore
empty, and Table~\ref{tab:s11} quantifies how far an implementation must go to
enter it: $s_{11}^*=2.7$--$3.6$ at $n=4096$, against $2.34$ measured.

\subsection{Attribution of the performance gap}\label{sec:attrib}

A factorization two orders of magnitude slower than a vendor sparse solver
invites the question of where the time goes. We ablated four factors on the same
device and the same circuit matrices, changing one at a time
(Table~\ref{tab:attrib}).

\begin{table}[t]
  \centering
  \small
  \caption{Multiplicative attribution of the gap to a vendor sparse solver.}
  \label{tab:attrib}
  \begin{tabular}{llrl}
    \toprule
    factor & measured ratio & evidence & nature \\
    \midrule
    A ordering / fill   & $\le 5.9$ (median $1.0$) & blocking vs global AMD, 13 matrices & algorithm \\
    B blocking          & $1.5$ (FP32 only)        & TILE $8/16/32$, same kernel          & implementation \\
    C pivoting          & $1.0$                    & partial vs static, 6 matrices        & none \\
    D kernel efficiency & $6.9$ (residual)         & $2.82$ of $29.03$\,TFLOP/s           & implementation \\
    \midrule
    product, typical & $\approx 10$ & $A=1$, $C=1$ & \\
    product, worst   & $\approx 61$ & \texttt{rajat25}, $A=5.9$ & \\
    \bottomrule
  \end{tabular}
\end{table}

Two conclusions follow. First, the algorithmic levers are nearly exhausted:
ordering changes the fill by at most $5.9\times$ and usually by nothing, and a
static pivot strategy buys stability at a $3.9\times$ fill cost on the single
matrix where it matters---a trade the growth factor does not require. Second,
the residual gap is kernel engineering, not hardware: the achieved rate is
$9.3\%$ of the vendor GEMM on the same device, and the roofline places every
matrix two orders of magnitude to the right of the ridge point, so the compute
roof is far above where the kernel operates.

\subsection{Threats to validity}

\paragraph{Emulator.}
Rounding is applied per operation in double, so a genuine FMA would differ in
the last bit; and the emulator has no notion of overflow or subnormals, which
matters for narrow-exponent formats such as FP16 but not for the TF32-like
11-bit model used here. Conclusions about \emph{scaling} are robust; conclusions
about constants are not.

\paragraph{Scale and parallelism.}
Dense experiments stop at $n=576$; nothing here models communication, memory
traffic, or the pipeline bubbles of a precision switch. The cost model is
flop-based, and Table~\ref{tab:platforms} is the only wall-clock measurement.

\paragraph{The iteration count $k(\kappa)$.}
Proposition~\ref{prop:corner} takes $k$ as input, and we supply it from
measurement rather than from a closed-form bound; the available bounds
\cite{hm2022} constrain feasibility, not the iteration count. The crossover
$n^*$ should therefore be read as accurate to the accuracy of the measured $k$.

\paragraph{Unmeasured upper end of $s_{11}$.}
The tensor-core path is not reachable through the documented API on our stack,
so $s_{11}$ is bounded from below by $2.34$ and from above only by vendor
specifications. The threshold in Table~\ref{tab:s11} does not depend on this
gap, but the position of the low-precision region in Figure~\ref{fig:phase} does.


\section{Application: when is FP32 safe?}

With $\rho\approx1$ on ten of thirteen matrices, $\rho u_{32}\approx6\times10^{-8}$,
two orders of magnitude below the $10^{-6}$ residual target; only the three
near-singular matrices require FP64. \todo{corroborate with end-to-end timings}

\section{Related work and conclusion}

\paragraph{Mixed-precision iterative refinement.}
The three-precision analysis of Carson and Higham \cite{ch2018} separates the
factorization, working and residual precisions and establishes the threshold
$\kappa u_f\ll1$ for stationary refinement; the five-precision framework of
Higham and Mary \cite{hm2022} adds the precisions of the GMRES operator and of
the preconditioner application, and gives the condition
$\kappa^2u_f^2(u_g+\kappa u_p)\ll1$ for GMRES-based refinement. Both take the
factors as given. Our result is orthogonal: we hold the algorithm fixed and vary
the construction of the factors, showing that the conditions above are met by
structurally different objects and that neither substitutes for the other.

\paragraph{Emulated high precision.}
Ozaki's error-free transformation \cite{ozaki2012} and its tensor-core
realizations evaluate a high-precision product as several low-precision ones,
provably and without a residual. This is complementary to refinement: emulation
attacks the $O(n^3)$ term by making each product exact, whereas we attack it by
making each product cheap and recovering accuracy from the residual. The choice
between the two is a hardware question---emulation needs no FP64 unit, refinement
does---and the cost model of Section~\ref{sec:structure} makes it explicit.

\paragraph{Precision assignment.}
Three threads assign precision \emph{spatially}: tile-centric mixed-precision
Cholesky in the Gordon Bell campaigns, and most recently the error-controlled
ladder of Abdurakhmanov et al. \cite{mx2026}, which selects, per tile, the
smallest MX format that satisfies a rounding-and-underflow budget. That work is
the closest to ours and complementary: they decide \emph{which tiles} get which
precision, we decide \emph{how the factors at a given precision are
constructed}. A natural composition is to use their ladder to choose the
per-tile precision and our dichotomy to choose the construction.

\paragraph{Deployment.}
At scale, mixed-precision solves are already the default in HPL-MxP
\cite{hplmxp}, which factors in half precision with FP32 accumulation and refines
in FP64 with GMRES; our analysis explains why that combination, and not the
storage-only variant, is the one that works. In application codes the gains are
more modest and more uneven, as the surveys of \cite{kashi2024} document.

\paragraph{Conclusion.}
Three statements summarize the paper. First, factor precision is not a scalar:
compressed and recomputed factors at the same unit roundoff have the same
backward error and opposite behaviour, and the mechanism is a pivot floor.
Second, the attainable residual is governed by the componentwise growth factor
$\rho$, computable from the factors, which predicts the residual to within two
orders of magnitude across eleven orders of magnitude in $\rho$---and on circuit
matrices it is $\rho$, not $\kappa$, that separates the ten tractable matrices
from the three hard ones. Third, the optimal precision is a function of the pair
$(\kappa,s)$ and, at fixed hardware, of the \emph{implementation} through the
threshold $s_{11}^*$: lowering precision pays only when the low-precision kernel
is genuinely faster, which is a property of the code rather than of the device.

\begin{thebibliography}{9}
\bibitem{ch2018} E. Carson, N. J. Higham. Accelerating the solution of linear systems
by iterative refinement in three precisions. \emph{SIAM J. Sci. Comput.},
40(2):A817--A847, 2018.
\bibitem{amestoy2024} P. Amestoy et al. Five-precision GMRES-based iterative
refinement. \emph{SIAM J. Matrix Anal. Appl.}, 45(1):529--552, 2024.
\bibitem{hm2022} N. J. Higham, T. Mary. Mixed precision algorithms in numerical linear
algebra. \emph{Acta Numerica}, 31:347--414, 2022.
\bibitem{ozaki2012} K. Ozaki et al. Error-free transformations of matrix
multiplication. \emph{Numer. Algorithms}, 59(1):95--118, 2012.
\bibitem{mx2026} A. Abdurakhmanov et al. Error-controlled mixed-precision Cholesky
factorization with MX microscaling formats. \emph{ACM TOMS}, 2026.
\bibitem{hplmxp} J. Dongarra, P. Luszczek. HPL-MxP benchmark: mixed-precision
algorithms, iterative refinement, and scalable data generation. \emph{IJHPCA}, 2025.
\bibitem{kashi2024} A. Kashi et al. Mixed-precision numerics in scientific
applications: survey and perspectives. \emph{J. Supercomput.}, 2025.
\end{thebibliography}

\end{document}
